%
En vannkanne er fylt med vann som står i ro. Punkt A ligger inne i kanna, og punkt B ligger i tuten, i samme høyde som A (se figuren).

\begin{center}
\begin{tikzpicture}[scale=0.9, font=\small]
  % Vann (klippes ved vannflaten y = 2.4)
  \begin{scope}
    \clip (-1,-1) rectangle (5,2.4);
    \fill[oppgcol!18] (0,0) rectangle (2.4,3.0);
    \fill[oppgcol!18] (2.35,0.3) -- (3.75,3.0) -- (3.15,3.0) -- (2.35,1.5) -- cycle;
  \end{scope}
  \draw[oppgcol!70!black, thick] (0,2.4) -- (2.4,2.4);
  \draw[oppgcol!70!black, thick] (2.85,2.4) -- (3.45,2.4);
  % Kanne
  \draw[very thick, black!75] (0,3.0) -- (0,0) -- (2.4,0) -- (2.4,0.3) -- (3.75,3.0);
  \draw[very thick, black!75] (2.4,3.0) -- (2.4,1.5) -- (3.15,3.0);
  % Hank
  \draw[very thick, black!75] (0,2.6) .. controls (-0.9,2.6) and (-0.9,0.8) .. (0,0.8);
  % Punkter
  \draw[dashed, gray!70] (0.6,1.8) -- (2.85,1.8);
  \fill[black] (1.0,1.8) circle (2pt) node[below] {A};
  \fill[black] (2.85,1.8) circle (2pt) node[right=4pt] {B};
\end{tikzpicture}
\end{center}

\begin{enumerate}[a)]
    \item En trykkmåler viser at trykket i A er \SI{1.05e5}{\pascal}. Hva er trykket i B? Begrunn svaret.
    \item Bruk svaret i a) til å forklare hvorfor vannet i tuten står like høyt som vannet i kanna.
    \item Du vipper kanna slik at vannet renner ut av tuten. Kan du fortsatt være sikker på at trykket i A og B er likt? Begrunn svaret.
\end{enumerate}
%
